\documentclass[12pt]{article}
\usepackage[utf8]{inputenc}
\usepackage[T1]{fontenc}
\usepackage{lmodern}
\usepackage{amsmath,amssymb}
\usepackage{graphicx}
\usepackage{hyperref}
\usepackage{microtype}
\usepackage{geometry}
\geometry{margin=2.5cm}
\title{Extracted PDF to LaTeX}
\date{}
\begin{document}
\maketitle
\tableofcontents
\newpage
% --- Page 1 ---
\section*{Page 1}
I'\textasciitilde{} THE BOA Rn Rf Sf ARCf! STUDIES Ar" \textbackslash{}cdot \textbackslash{}cdot f(lr? ; Pb. D. ,, . ffB '\textbackslash{}cdotII..' ' . < ' 1966 \textasciitilde{} \textbackslash{}cdot- - - - - - - - - -\_\_j 01 l l l \textbackslash{}cdot \textbackslash{}cdot ONIVi:Rc LB l CAMUKliJ<JE \textbackslash{}cdot te .L e re o h . ivc .\textbackslash{}cdot \textbackslash{}cdott o\textasciitilde{} bri o

\begin{figure}[h]
\centering
\includegraphics[width=0.9\linewidth]{images/page001_img01.png}
\caption{Extracted image from page 1}
\end{figure}

% --- Page 2 ---
\section*{Page 2}
om lica ion und c u c o oft i n\textasciitilde{}e niv r c r 't l ....\_,i d . In vh >t r 1 it is o n t t t i er u.te r ve ifric l for t oyle i\textasciitilde{}o \_ li\textbackslash{}cdot r r .... o 0 Vl\textbackslash{}cdot . ion . ter e l i th \textbackslash{}cdott tion of \textbackslash{}cdot n OU i ot o ic u i\textbackslash{}cdot ver \textbackslash{}cdot c t l\textbackslash{}cdot i ..., c 0 or r t, 0 t .L'O \textbackslash{}cdot r ... ur h re initi l l . he nro i\textbackslash{}cdot on tion of r \textbackslash{}cdot Vl.- l .... o i VC.Jti i\textbackslash{}cdot n i.; roxj atio . . . .... In Jt1 ptcr r\textbackslash{}cdot vit tion 1 r \textbackslash{}cdot n n r univ r 0 "' l. i inc ;t 0 0 o ic 1\textbackslash{}cdot 0 t l' - \textbackslash{}cdot 0 t viour n t c s to Lie u e i d . l Lt l t i o .... urrcnc o \textbackslash{}cdot i\textbackslash{}cdot n co olo c l o el . t i ... o n t u t i l. i i it 1 n ovi e t err. i v r l c n r \textbackslash{}cdot l. \textbackslash{}cdot

\begin{figure}[h]
\centering
\includegraphics[width=0.9\linewidth]{images/page002_img01.png}
\caption{Extracted image from page 2}
\end{figure}

% --- Page 3 ---
\section*{Page 3}
- J.. ... 0 0 i 1 i 0 t \textbackslash{}cdot lt"\} l l l it i t \textbackslash{}cdot \textbackslash{}cdot - \textbackslash{}cdot ' - ' \textbackslash{}cdot ' - l \textbackslash{}cdot v i ' - - \textbackslash{}cdot \textbackslash{}cdot \textbackslash{}cdot - \textbackslash{}cdot - t t

\begin{figure}[h]
\centering
\includegraphics[width=0.9\linewidth]{images/page003_img01.png}
\caption{Extracted image from page 3}
\end{figure}

% --- Page 4 ---
\section*{Page 4}
i \textbackslash{}cdot ' 0 ' - \textbackslash{}cdot \textbackslash{}cdot ' \textbackslash{}cdot \textbackslash{}cdot . - I - ' .l \textbackslash{}cdot \textbackslash{}cdot \textbackslash{}cdot \textbackslash{}cdot i 0 '

\begin{figure}[h]
\centering
\includegraphics[width=0.9\linewidth]{images/page004_img01.png}
\caption{Extracted image from page 4}
\end{figure}

% --- Page 5 ---
\section*{Page 5}
. .. ' \textbackslash{}cdot - 1 \textbackslash{}cdot - \textbackslash{}cdot - ' \textbackslash{}cdot .. t \textbackslash{}cdot - - k :;::. I 0

\begin{figure}[h]
\centering
\includegraphics[width=0.9\linewidth]{images/page005_img01.png}
\caption{Extracted image from page 5}
\end{figure}

% --- Page 6 ---
\section*{Page 6}
- \textbackslash{}cdot \textbackslash{}cdot l \textbackslash{}cdot !\textbackslash{}" ::. \textbackslash{}cdot 0 l \textbackslash{}cdot \textbackslash{}cdot \textbackslash{}cdot 4 - - \textbackslash{}cdot \textbackslash{}cdot - 1. \textbackslash{}cdot \textbackslash{}cdot i \textbackslash{}cdot - i \textbackslash{}cdot . ' 1 \textbackslash{}cdot ' i \textbackslash{}cdot l 0 i ' \textbackslash{}cdot

\begin{figure}[h]
\centering
\includegraphics[width=0.9\linewidth]{images/page006_img01.png}
\caption{Extracted image from page 6}
\end{figure}

% --- Page 7 ---
\section*{Page 7}
.u ..... 1.C () ..Lr \textbackslash{}cdot l i t thank U"")ervisor , .vr. \textbackslash{}cdot a \textbackslash{}cdot \textbackslash{}cdot CJ. f r n lp n.n urin oy of r seurcn. i \textasciitilde{}dvice pc\textasciitilde{}io I oulu like to thm r . \textbackslash{}cdot vurter ond \textasciitilde{}lso .Jr . \textbackslash{}cdot \textbackslash{}cdot \textbackslash{}cdot \textasciitilde{} \textbackslash{}cdot Sl.1i\textasciitilde{} for \_..\_:ny u\textasciitilde{} ful l\textbackslash{}cdot. SCUG i\textbackslash{}cdot ons . I i ndebtcu o R. G. cLcna or the cr..lcul t i on of the innchi i entities in vh ptcr j . 1.nc re. curer\_ ue...>cri bed i i thi .... t lP. is t.'t n carri on l.ilc I h'"'l .... utu ent h: rom \textbackslash{}l 0 e\textasciitilde{}rch t 10 J . u . l . \{ . \textbackslash{}cdot S. LJ , kM13 H\textasciitilde{} c c 19 1 :;t h t 0 :::> \textbackslash{}cdot

\begin{figure}[h]
\centering
\includegraphics[width=0.9\linewidth]{images/page007_img01.png}
\caption{Extracted image from page 7}
\end{figure}

% --- Page 8 ---
\section*{Page 8}
1 CI-i.A.PT\textasciitilde{}rl The Hoyle - Narlikar Theo\textasciitilde{}y of Gravitation 1. Introduction The success of Maxwell's equations has led to electrodynamics being normally formulated in terms of fields that have decrees of freedom independent of the in \textasciitilde{}articles them . However , Gauss suggested tnat an ac\textasciitilde{} ion-at-a-distance t :1eory in \textbackslash{}'Jhich the action travelled at a finite velocity might b e possible . This i dea was developed by 1heeler and 1 1 2 Feynman ( , ) who derived their theory from an action-principle tha\textbackslash{}cdott involved only direct int eractions between pairs of par\textbackslash{}cdotticles . feature of this theory was that the ' pseudo '-fields A int roduced are the half- retarded plus half-advanced fields claculated from bhe worl d - lines of the particles . However, (3) \textbackslash{}cdot\textbackslash{}cdot! :1eeler and JJ'eynman , and , in a different v-1a;; , Ho,-artli were able to show that, provided certain c osmolo3ical condit ions were satisfied, these field s could combine to 4 6ive the observed field. Hoyle and Narlikar ( ) extended \textbackslash{}cdottt1e ory to general space - t;imes and obtained similar theories for thE:ir ' C'-field (5)and for the gravi\textbackslash{}cdottational f:fueld (G). It is with these theories that this chapter is concerned. - I - ...... . ..... I I ,I\textbackslash{} I< y L J A\textasciitilde{}V "A"'l / -

\begin{figure}[h]
\centering
\includegraphics[width=0.9\linewidth]{images/page008_img01.png}
\caption{Extracted image from page 8}
\end{figure}

% --- Page 9 ---
\section*{Page 9}
It \textbackslash{}'Jill bG .. .;ho.:n that in an exp;.11ding universe t\textbackslash{}cdot!1e advanced fields are infinite, and the retarded fields finite. This is because, unlike electric charges, all Toasses tho h[tVA same ...:;\textbackslash{}cdot i.;;;n . 2. The Boundary Condition rloyle and Narlikar derive their theory from the action: \textbackslash{}cdot wnere the integration is over the world-lines of oarticles b. a. In this expression \textasciitilde{} is a Green function . \textbackslash{}cdot \textbackslash{}cdot \textbackslash{}cdot I that satisfies \textbackslash{}cdotthe wave equation: l(( . /) x £ x. - - /-3 3 j wl1ere is the dei:;erminant of .j \textbackslash{}cdot -=>inc:e double sum t\textasciitilde{}e A in the action is symcetrical between all pairs of 4\textbackslash{}cdot(0...b) pi:irticles a.,b , or1ly thc..t part of that is b symmetrical betv.reen c.... and v1ill contribute to the action i.e. the action can be writ\textbackslash{}cdotten A b) \_ \textbackslash{}vhere \textasciitilde{}\textbackslash{}cdot\textbackslash{}cdot)!(-( \textasciitilde{} \textbackslash{}cdotlt 'rhus G must be the time- symmetric Green function, and can 1 ()"' ..!.. q q be i.,rrittien: '1 :: \textasciitilde{} re-t' -r o..dv v1here

\begin{figure}[h]
\centering
\includegraphics[width=0.9\linewidth]{images/page009_img01.png}
\caption{Extracted image from page 9}
\end{figure}

% --- Page 10 ---
\section*{Page 10}
are the r3tarded adval1ced Green f\textbackslash{}cdotL1nctions. r\textasciitilde{}nd Ey that the action be stationary under variations requirin\textasciitilde{} , of the Hoyle and Narlikar obtain the field - equations: g\textasciitilde{} i t [ £. i g"K R) x) x )] (r<il< - \{Vl,lo.)( M(b) ( -:i: b 0.. - < <: (a.)\{ \_!\_ [ 9 ' (b) r (b) ) (a) ( b) G) ( ::: - ft,K -r \textasciitilde{} c\_ 3 - /'II. fVl, - f'rL. \textbackslash{}cdott a(. fYl, /l1. . o..fb J vk. \} r ;\textasciitilde{}K )i, ; r\textbackslash{} - (o.);r (b)J m M., > , . Ho\textbackslash{}cdot,;ever, a.s a w11ere consequence of the particular choice of Green function, the contraction of the field-equations is satisfied identicall : . '1.'here a ::c thus only 9 equations for the 10 comnonents of Cr \textbackslash{}cdot '. ,J \textasciitilde{}J and t;hc S,',rstem is indeterminate . \{o..) Z : .. oyle and Narlikar therefore im')OSe rrt. = fvt =cor:st., 0 ,As the tenth equation. By then the 'smooth- fluid' m\textasciitilde{}kin < ( . t. th t . b tt \textbackslash{}cdot L o.) (b) 2 m approxima ion' a l.S y pu ing L L\_ (V\textbackslash{}, JI.- /11.. a...:tb CJ) they obtain the Einstein field- equations: - \_.!\_R\_ .) . . \textasciitilde{} \textasciitilde{}\{R. gi-1'\textbackslash{} : \textbackslash{}cdot (, \{IA..O .\_v, l - / t, K There is an important difference , ho\textbackslash{}vevor, bet\textbackslash{}veen these field-equations in the direct - particle ir1teraction tc.eory and ii1 tl1e usual general tneory of re L\&.tivity . In the seneral tneory of relativity, any metric that satisfies t he

\begin{figure}[h]
\centering
\includegraphics[width=0.9\linewidth]{images/page010_img01.png}
\caption{Extracted image from page 10}
\end{figure}

% --- Page 11 ---
\section*{Page 11}
the field-equations is admissible, but in the tlirect -particl \textasciitilde{}\textasciitilde{} interaction theory only solutions of the field - equations t\textasciitilde{}ose are admissible that satisfy t he additional requirement : (o..) £ - (x) - /VL = .. - This requirement is highly restrictive; it will be shown it is not satisfied for the cosmological solutions of th\textasciitilde{}t the binstein field- equations, and it appears it th\textasciitilde{}1t c\textasciitilde{}nnot \textbackslash{}cdotbe satisfied for any mod.els of tl1e univer:;e tho.t eitl1er \textbackslash{}cdot contain an infinite amount of matte1\textbackslash{}cdot or undergo infinite expansi\textbackslash{}cdot on . The difficulty is similar to that in occurrin\textasciitilde{} :r-rev1tonian t\_icory \textbackslash{}vhen it is recognized that the universe might be infinite . The Newtonian p otential obeys equation: th\textasciitilde{} 1 : -- "f D is the density. '\textbackslash{}cdot"'he re '

\begin{figure}[h]
\centering
\includegraphics[width=0.9\linewidth]{images/page011_img01.png}
\caption{Extracted image from page 11}
\end{figure}

% --- Page 12 ---
\section*{Page 12}
. In an infinite staLic unive\textasciitilde{}se , wo1Jld be since infini\textasciitilde{}e, the source has the same difficulty was al\textasciitilde{}ays \textasciitilde{}ign . \textasciitilde{}he \textasciitilde{}ssolved v1hen it v1as reolized the univ ersP. \textbackslash{}cdot.-1as expa11d.in;;, since th<\textasciitilde{}t in .:tn exoanding universe the retarded solution of the above is finite by a sort of ' red- shift ' eq\textasciitilde{}ation ef\textasciitilde{}ect . \textasciitilde{}he advanced solution will be infinite by a ' blue- shift ' effect . 'l'his is unimportant in Ne\textbackslash{}vtonian theory, since one i:..; free to choose vhe solution of the equation and so ms.y icnore \textbackslash{}cdotthe infinite advanced solution and take simply the finite retarded solution . \textasciitilde{}iffiilarly in the direct - particle interaction tQeory the rn, - field satisfies the equation: f R. M N O m = +- N where is the density of \textbackslash{}\textbackslash{}cdot.;orld- lines of particles . As in the Newtonian case, one may expect that the effect of the ex1)a11sion of the universe will be to mc:1ke \textbackslash{}cdotthe retarded solution . finite and the advanced solu\textbackslash{}cdottion infinite . However, one is now not free to choose the finite retarded solution, for the equation is derived from a direct - particle interaction action principle symmetric between pairs of particles , and one must choose for fV\textbackslash{}.. half the sum of the retarded and advanced solutions. Vie would expect t \textbackslash{}cdotJ.is to be infinite , and this is shown to be so in the next section .

\begin{figure}[h]
\centering
\includegraphics[width=0.9\linewidth]{images/page012_img01.png}
\caption{Extracted image from page 12}
\end{figure}

% --- Page 13 ---
\section*{Page 13}
\[
'fhe Gosmolos;ical i.:),olu-cions j . 'l1he .rtobertson- \'/alker cosmologic<:'l.l 1netrics ho.ve the forin - . \cdot Since t.1ey are conformalJ.y flat, one caoose coordinates c~n in \vhich tr.Ley become :z i 2 :L J. . Q 2 cls - 7o..b ; J1-< dx^{\circ}'- cLx , t1ere io the flat - space metric tensor and \cdot..v - -- - t~\_(c) .c 7) / i tl~ I K ( r-f )iJ [I ~ ~ ii ( ~ - ( )~J 't \cdott For example, for the Sitter un\cdoti verse Einstein-d~ - c 2 R(l)-:: k ;: 0 3 / ~ T ( ) 'L G .;..:\_ "f. <. rD T I ~) \_\_ Q ("'C= JJt For the steady- state (de Sitter) un\cdoti verse ..t K :. 0/ Rt.t); . (-o0.<.t:"Lo0) e1 \cdot
\]
\begin{figure}[h]
\centering
\includegraphics[width=0.9\linewidth]{images/page013_img01.png}
\caption{Extracted image from page 13}
\end{figure}

% --- Page 14 ---
\section*{Page 14}
\[
\cdot- R \cdot- ( - o) -- -- I oO <. L. L\_ :--'!'- [, - 7) (: (i: Te - .- \cdot- - f F c;* (~,, b) obeys t he equation 'I.he Green function 1 t b) ~<t{Q = oc;-#('(o..,h)~ R~*(CA,b) ;-:\_\_ 3 t his it follows that ~rom I J2'-i dX dx.b JL-'iJ ~ 4 (a,b) 5 52.-t ..\_ , If we let q~; v :1e11 ' - - Thi s is s imply the flat-s-oace Gr -.. en function equation , a nc1 \cdotn e11c e f ) - Jl-'?c,) *( ii~J2--,- c.) 0 . r 't. 'l '1 ,, ) l - J2. ( ti)f ~ If - g Cr) --c- vc, ) \cdot f' L - f J J2. (. "C ).\_ The 1 \cdotI\cdot fie l d is given by r J /V\. I ( dx: ~re,t\_ J\_ '~x:N/-3 = lh.0-dv: \cdot=- \t \cdotr the Einstein-de Bitter without creation (e.g. lt'or uni verses N 3 ~ I~ \cdot- -- ,., /1., ,,. const . .e or universe),
\]
\begin{figure}[h]
\centering
\includegraphics[width=0.9\linewidth]{images/page014_img01.png}
\caption{Extracted image from page 14}
\end{figure}

% --- Page 15 ---
\section*{Page 15}
con.:;;t. W\textbackslash{}Gre the integr ation over the fu t ure light cone . This \textasciitilde{}s 2-: . normally be infinite in an expandir13 :.niv erse , e in \textasciitilde{}\textbackslash{}cdot1ill \textasciitilde{} 0itter universe. th\textasciitilde{} \textasciitilde{}instein-de -2 oO vc ) cL 1 ( f - M. o...).; ( '-f,) -- -'1: \textasciitilde{} .2 I l I I " ... /., I - - In the s't:;eady- state un\textbackslash{}cdoti verse 3 \textbackslash{}cdot- - -l 0 - I ( ) cL c:.z . \textbackslash{}cdot- I - n. l , - ,l/\textbackslash{}- - ( '1;( ) J ./. /ll\textbackslash{}... 1. CL T. l -r. I \textasciitilde{} - I - \textbackslash{}cdot .B y on the other hand , \textbackslash{}ve have c\textasciitilde{}o»ntrast, where the integration is over the past light cone. This will

\begin{figure}[h]
\centering
\includegraphics[width=0.9\linewidth]{images/page015_img01.png}
\caption{Extracted image from page 15}
\end{figure}

% --- Page 16 ---
\section*{Page 16}
normally be finite , e . g . in the univers= \textasciitilde{}in\textasciitilde{}tein-de 0i\textasciitilde{}cer -2 'L . f. ( -Q. t'.2.. - 'L ) c\{ '"'L ::. ' I ' '\_ /?\_, -(L 1. J. T I I 0 while in the steady- state universe - \textbackslash{}cdotr - -1 - IV1.. ( 't, ) -- I I n. I \{; ;\textasciitilde{} '"f \textasciitilde{} \textbackslash{}cdot-cO I l;i. -i \textbackslash{}cdot-:i. - - J 2.. (L Thus i t cun be seen that the solution consi::; . of ll'\textbackslash{}.. =- th0 equation is not , in a cosmological metric , the half-advanced plus half- retar\textbackslash{}cdotded solution since this would be infinite . In fact, in the case of the Einntein- de and steady-state metrics, \textasciitilde{}it\textbackslash{}cdotter it is the pure retarded solution . 4 . 'C '-b'ield \textasciitilde{}he Hoyle and ]arlikar derive their direct - particle interaction theory of the 'C' - field from the action

\begin{figure}[h]
\centering
\includegraphics[width=0.9\linewidth]{images/page016_img01.png}
\caption{Extracted image from page 16}
\end{figure}

% --- Page 17 ---
\section*{Page 17}
\[
h where tae suffixes refer to diffqrentiation of Q \cdot'^{\circ}' I (o..,b) 6 ~~ on the world- lines of rcGoectivel~' \cdot l'-1.. 1 G......... is a Gr een function obeying the equation DG(X,X') ,~e define tl1e ' C' - field by \cdot z G'""' (x)^{\circ}-) clQ & ) C(x) ,~ = 0.... .I and the by mat~er-current a J"'<~) \cdot z\_ 4 L ~, b)cLb \cdot~, S' e( 3 1) J j - c &-J lj ) cl\_ .rhen X, '< ( J K 'X 4J 1 o o c - I'\ - J - JK :~e tl1us .:>ee that the sources of the 1 C' - field. are the plc.ces matter is created or destroyed . w~ere i1.s in the case of the '-field , the Green func\cdott; i on ' M, must be timc- S;) mmetric , that is 7 q( - ^{\circ}', b) I - - ,
\]
\begin{figure}[h]
\centering
\includegraphics[width=0.9\linewidth]{images/page017_img01.png}
\caption{Extracted image from page 17}
\end{figure}

% --- Page 18 ---
\section*{Page 18}
hoyle and Narl ikar claim thut if action of the th\textasciitilde{} ' C' - fielQ is included al vii th the action of tl1e ' '-field, on,\textasciitilde{} r\textasciitilde{} a univerue \textbackslash{}1ill be obtained that approximates to th8 3teadystate universe on a large scale althOUBh tncre may be local irrecularities. In this universe , the value of C will be f i:ni 1..ie and its gradient time - like and of unit ma5nit;ude . Given this universe, we may check it for consistency claculatin6 the and retarded ' C'-fields and advan\textasciitilde{}ed fincin\textasciitilde{} if their sum is finite. shall not do tl1is directly \textasciitilde{}e \textasciitilde{}ut will that the advanced field is infinite while the \textasciitilde{}how re\textasciitilde{}arded field is finite . Consj.der a region in space - time bounded by a three dirnensional space - like hypersurface ]) at the present time, P £.. und t(1e past light cone of some point to the futl1re j) 9f \textbackslash{}cdot By Gauss's the orem /-3 w dx C \textasciitilde{} \textasciitilde{} I/ t/ 1 Let the advanced field pro(luced by sources v1i thin be C \textbackslash{}cdot l C Then and will be zero on \textasciitilde{} , and hence

\begin{figure}[h]
\centering
\includegraphics[width=0.9\linewidth]{images/page018_img01.png}
\caption{Extracted image from page 18}
\end{figure}

% --- Page 19 ---
\section*{Page 19}
.-. 0 - 1-\textbackslash{} J .dut is the rate of creation of matter= ft,. (canst . ) i\textbackslash{}cdot n \}K the steady- state universe , and hence V. - rv - ? the point is taken further into the future , the volume 11.S of the re6ionV t\textasciitilde{}nds to infinity . Hov1ever, the area of tlJ.e hyp.ersurface J) tends to a finite limit owing to horizon c d I must be infinite. effects . the gradient \textasciitilde{}herefore dnA similar calculation shows the gradient of the retarded field to be finite . '.I.'heir sums cannot therefore :-\textbackslash{}cdotiv;e the neor;,r . field of unit gradient required by the I \textbackslash{}cdot Hoyle -t\textasciitilde{}arlilcar v It io worth noting that this result was obtained without assumptions of distribution of matter or a\textasciitilde{}smooth of conf flatness. ormul.\textasciitilde{}

\begin{figure}[h]
\centering
\includegraphics[width=0.9\linewidth]{images/page019_img01.png}
\caption{Extracted image from page 19}
\end{figure}

% --- Page 20 ---
\section*{Page 20}
5. Conclusion It iJ one of the weaknesses of the Einstein theory of relativity that \&lthough it furnishes field equations it does not provide bour1dary conditions for them. 'rhus it does not Give a unique model for the universe but allows a whole series of rnoa.el:::> . Clearly a theory th;:,t provided boundary "Gior"s conc...L and. t.1ut:> re.stricted the possible solutions \textbackslash{}..,rould l)e VP>r\textbackslash{}r y ., attractive. The Hoyle- Narlikar t11eory does just .. t(the th\textasciitilde{} I m J.\_ M ) . l."'equirement that = ;:- mt\textasciitilde{}t 't ). O...d ./ l.S equivalent to a boundary condition) . Unfortunutely, as .. e have seen above , this condition excludes th\&t tho\textasciitilde{}e ;\textasciitilde{} odels seem correspond to the actual universe , namely the 1,0 models. R obertso)1-\textasciitilde{}\textbackslash{}cdotJ al\textasciitilde{}::er . calculations given above have considered the Th\textasciitilde{} univct\textbackslash{}cdots\textasciitilde{} as being .f'illP.d \textbackslash{}\textbackslash{}cdot1ith a 11niform distribution of m<:.:.tter. '-."his is if we are able to make the 'smooth-fluid ' le\textasciitilde{}itimate approximation to obtain the equations. \textasciitilde{}instein \textasciitilde{}lternativel-- if this approximation is invalid, it cannot be said that the tneory yields the Ei nstein equations. It mi.5ht possibly be that local irre\_;ularities co11ld :nal\{e finite , but this has certainly not been demonstrated 'Yta..dV 1 and seems unlikely in vievJ of the fc..\textbackslash{}ct th1 ... t , in t11e Tio\textasciitilde{}rle - Narlikar direct - particle interaction theory of their 'C 1 -field,

\begin{figure}[h]
\centering
\includegraphics[width=0.9\linewidth]{images/page020_img01.png}
\caption{Extracted image from page 20}
\end{figure}

% --- Page 21 ---
\section*{Page 21}
which is derived from a very similar action- principle, it can be shown without assuming a smooth distribution that the advanced C ' field \textbackslash{}-.rill be infinite in an expandin3 univ erse 1 'IJi th crea\textbackslash{}cdott ion . \textasciitilde{}he r eas on that it is possible to f ormulate a dire ctparticle interaction theory of e lectrodynaffiics that does \textasciitilde{}ot mf encounter this difficulty havin3 the advcincea solution infinite is that in electrodyncJJics tl1ere are eo.ual n11mbers of so\textbackslash{}1rces of nositive and negative sie;n . Their fi.elds cu.r c\textasciitilde{}ncel each other out and the total field can be zero apart fron1 local irregularities . This sug:est that 1)ossible \textbackslash{}cdot:Jay i:.. to save the Hoyle - Narlikar theory would be to allow masses 0\textbackslash{}cdot1\textasciitilde{}.\textbackslash{}cdot The action would be s i\textbackslash{}cdot gn . both p osit ive and negative G b) clo. cL b \textasciitilde{}\textasciitilde{}, are gravitational char.\textasciitilde{}es a\textasciitilde{}alo\textasciitilde{}ous to q,. electric char \textbackslash{}cdotes . Particles of positive in a r>osit i ve - in a ne .ra\textbackslash{}cdott i ve 1 I rr1 fi1 ' -field und particles of negut ive Cj\_ ....., 1 fie ld W()\textbackslash{}lld h\&ve the noi... mal \textasciitilde{}ravi tational properti\textasciitilde{}s, is, t;t'i.:.,:.t \textbackslash{}cdott11e;f would have positive gravitational and inertial masses. '

\begin{figure}[h]
\centering
\includegraphics[width=0.9\linewidth]{images/page021_img01.png}
\caption{Extracted image from page 21}
\end{figure}

% --- Page 22 ---
\section*{Page 22}
1 A particle of in a positive '-field ne\textasciitilde{} ative ' M \textbackslash{}cdot:J(;u ld s\textbackslash{}cdottill J'ol lov1 a geodesi c. Therefore it wo11ld be 0.tt;racted :)y a particle of positive Its own gravi tat:;ional effect \textbackslash{}cdot howeve r would be to repel all other particles . Thus it would hav\textasciitilde{} the properties of the negative mas J described by Bondi(e) tt1at i s , negative gravitational 1110.so and n egat ive inertial mass . 3ince tnere does not seem to be any ma tter having toese propert ies in our region of space ( 1ere M coi:.st . / 0 ) v11 JL there must clea rly be separat i on on a very larc e scale . '-' .l. ,, would not be possible to identify part i cles of ne\textasciitilde{}\textasciitilde{}ltive with ant imat t er , since i t is known that antimat\textbackslash{}cdott er has positive inertial mass . liwev er , the i ntroduction of nesat ive masses would prob ably raise more difficulties than it wuld solve . '

\begin{figure}[h]
\centering
\includegraphics[width=0.9\linewidth]{images/page022_img01.png}
\caption{Extracted image from page 22}
\end{figure}

% --- Page 23 ---
\section*{Page 23}
REFERE1\textasciitilde{}CES 1 . J . A. and Rev . Mod . Phys . 17 157 1945 \textasciitilde{}ilieeler R. P . Feynman 2 . J . A. \textbackslash{}Jheeler and Rev . Mod . Phys \textbackslash{}cdot 21 425 1949 1 tl . P . :E' cynman . J . E. Ho arth Proc . Roy . ooc . A 267 365 1962 ./ 6 4 . F . Ho.1le and Proc . :::<oy \textbackslash{}cdot ...,oc . A 277 1 1964a J . V. Narlikar 5 . F . tloyle and 282 178 19\&+b Proc . \textasciitilde{}oy . Soc . A J . V. Narlikar 6 . F . .. ioyle and Proc . Roy . boc . A 282 191 1964c J.V . Narlikar 7 . L. Infeld and Phys . Rev . 68 250 1945 A. Schild 8 . H. Bondi Rev. Mod . Phys . 29 423 1957

\begin{figure}[h]
\centering
\includegraphics[width=0.9\linewidth]{images/page023_img01.png}
\caption{Extracted image from page 23}
\end{figure}

% --- Page 24 ---
\section*{Page 24}
CH\textbackslash{}P':'ER 2 PERTURBAT IONS -- Introduc.-t.i-o-n- 1 \textbackslash{}cdot Perturbations of a spatially isotropic and homogeneous expanding universe have been investigated in a Newtonic'Jl approximation by 1 2 Bonnor( ) and relativistically by Lifshitz( ), Liftshitz and 4 Khalatnikov(3) and Irvine( ). Their method was to consider small variations of the metric tensor. This has the disadvantage that the metric tensor is not a physically significant quantity, since one cannot dire ctly measure it, but only its second derivatives. It is thus not obvious what the physical interpretation of a given perturbation of the metric is . Indeed it need have no physical significance at all, but merely correspond to a coordinate trans formation. Instead it seerns preferable to deal in terms of' perturbations of the physically significant quantity, the curvature . 2. Notation Space-time is represented as a four- dimensional Riemannian space with metric tensor gab of signature +2 . Covariant differentiation in this space is indicated by a semi - colon. Square brackets around indices indicate antisymmetrisation and round brackets symmetrisation. The conventions for the Riemann and Ricci tensors are: - 1/91.b cet is the alternating tensor. Units are such that k the gravitational constant and c, the speed of l ight are one . Ul'\textbackslash{}I VtlC>' I Y Ll.:il'/\textbackslash{}cdot\textasciitilde{}Y CAMBRIDGE

\begin{figure}[h]
\centering
\includegraphics[width=0.9\linewidth]{images/page024_img01.png}
\caption{Extracted image from page 24}
\end{figure}

% --- Page 25 ---
\section*{Page 25}
3. The Field\_.\textasciitilde{}gua\textasciitilde{}.io\_ns We assume the Einstein equations: where Tab is the energy momentum tensor of matter. We will assume that the matter consists of a perfect fluid. Then, a - where Ua is the of the fluid, U U - - 1 : vel\textasciitilde{}city a is the density . \textasciitilde{} ft is the pressure h is the projection operator G\textbackslash{}b : Sob Ua.. l..tb t into the hyperplane orthogonal t o U \textbackslash{}cdot a . uh hAb ::0, We decompose the gradient of the velocity vector Ua as \textbackslash{}cdot t + 8 - UQ.;b:. Wob-+ <5'"o.b h\textasciitilde{}b Uc.. Ub where is the acceleration, is the expansion, hJ hC. I is the shear, <r'o. b \textasciitilde{} IA ( c j <l) d. b - -S h:. is the rotation of the wo.b \textasciitilde{} l1t u [c;.d] flow lines U \textbackslash{}cdot We define the. rotation vector Wo. as a c.ol b I Wo. -:;. -;: 7\}o.bc.cl W Ll \textbackslash{}cdot We may decompose the Riemann tensor R into the Ricci tensor abcd Rab E1..nd t11e \textbackslash{}cdot:.'cyl tensor Ca bcd : R cbco( = Ca.\textasciitilde{}ccl - 9q(ct Rc.J I:;, - \textasciitilde{} b[c l\textasciitilde{}cl]\textasciitilde{} - R;3 'JG\textbackslash{}.[,5olJb ) Cabe.et\textasciitilde{} C[.o.b)\textasciitilde{}d.J 7 u C \textbackslash{}cdot. b'o. '==' 0 -::. Co.[ bed] \textbackslash{}cdot

\begin{figure}[h]
\centering
\includegraphics[width=0.9\linewidth]{images/page025_img01.png}
\caption{Extracted image from page 25}
\end{figure}

% --- Page 26 ---
\section*{Page 26}
Cabcd is that part of the curvature that is not determined locally by the matter. It may thus be taken as representing the free gravit ational field (Jordan, Ehlers and Kundt ( 5)). 'Ne may decompose it into its "electric" and "magnetic11 components. ' u t cd \_ tc.-\textbackslash{}cdot .. ,tJ cC'- Ed] C 8 [o.. b] \textbackslash{}.A - 4- O(o.. !,) o.b - ' ') , ) E H = o -=-- a. o. Q... (). ) each have five independent components. E ab We regard the Bianchi identities, as field equations for the free gravitational field. Then 6 (Kundt and Trfimper,C )).

\begin{figure}[h]
\centering
\includegraphics[width=0.9\linewidth]{images/page026_img01.png}
\caption{Extracted image from page 26}
\end{figure}

% --- Page 27 ---
\section*{Page 27}
\[
Using the decompositions given above, we may write these in a form analogous to the Maxwell equations. h b E. hcd H b '? b c Lide \ b l ( 1 ) 1. + ~ o.h W - o..~c.d U <Y e n = '5" n bc;ol l'A }"-;" J (j. (2) ' E - E Y/ uP a- e.r YJ c:l'f c(o. gb)c - o c.c:le bp~r Uc \cdot - H d Uc u\cdot e -- t(fi (3) + 2. YJ bcc=I. e ~) () ob Q 'J \cdot E (o. L H c:o..b - h (o. .f ') 'b) c d. e \). c. -f o~e + Hott. 8 ..- H w b)c. H ... .. H - "1 u <Y c. (P- o-b)c. p.r; cl e Y) b p q ,. 1,.4. c. p cl ~ ~ (4) \cdot .. 2H("J '". - - t- o. ~ he. ,,I e U. tA e 0 \cdot where .1. indicates projection by hab orthogonal to Ua. (7)) ( c.f. Trfunper, \cdot The contracted B\cdotianchi identities give , \cdotb T - cob' -= 0 ) fa+ (f+~)8 =. 0 (5) > r.; (f.+ h ~) Gl + 'p b a :: 0 (6) 0 The definitien of the Riemann tensor is, = 'R U o.;[bc.] 2. f' be. uP Ol \cdot Using the decompositions as above we may obtain what may be regarded as equations of motion" ,- 11
\]
\begin{figure}[h]
\centering
\includegraphics[width=0.9\linewidth]{images/page027_img01.png}
\caption{Extracted image from page 27}
\end{figure}

% --- Page 28 ---
\section*{Page 28}
\[
e \cdot = 2. '2. I n2. zw -2a- -3l7 + (7) ) (8) ' \cdot. .. .. (9) ~ \cdot' ' 2. 06 \Vhere 2 W =: W ab W ' \/Ve also obtain \\rr1at may be regarded as eq\_uati ons of constraint. ( 1 0) ( 11 ) ) n L-t c.( d;e = - n (o \cdot1b)c.cJ.e u \.. wf \cdot ( 1 2) consider perturbations or a universe that in the undisturbed state ~e in conformally flat, that is 0 0 \cdot abcd (3), By equations (1) - this implies, \cdot
\]
\begin{figure}[h]
\centering
\includegraphics[width=0.9\linewidth]{images/page028_img01.png}
\caption{Extracted image from page 28}
\end{figure}

% --- Page 29 ---
\section*{Page 29}
If \textbackslash{}Ve assume an equation of state of the form, \textasciitilde{} it"\textbackslash{}. ( \textasciitilde{}) , e. = \textbackslash{}cdot 1:Jien by ( 6) , ( 1 0) , :::. o .. Uo. This implies that the universe is spatially homogeneous and isotropic since there is no direction defined in the 3- space orthogonal to Ua. In this universe we consider small perturbations of the motion of tl1e fluid and of the '.ifeyl tensore Ne neglect products of small 1 quantities and perform derivatives with respect to the undisturbed metric. Since all the quantities we are interested in with the exception of the scalars, µ, \textasciitilde{}' e have unperturbed value zero, we avoid perturbations that merely represent coordinate transformation and have no physical significance. (4) To the first order the equations (1) - and (7) - (9) are h b - b E j -I s. ( 13) f;b ' ci-b - 0 ( 14) ) \textbackslash{}cdot r5) ( 1 b+E L8t- ) - c;a.. 0 r,:> 6) ( 1 ' e2. \textbackslash{}cdot - .. e \textbackslash{}cdot \textbackslash{}cdot Q. - -I + lA. .I ( 17) 3 Q - - - 'l. 8 Ll0-b + 8) ( 1 - 3 J ( 19)

\begin{figure}[h]
\centering
\includegraphics[width=0.9\linewidth]{images/page029_img01.png}
\caption{Extracted image from page 29}
\end{figure}

% --- Page 30 ---
\section*{Page 30}
From these we see. that perturbations of rotation or of Eab or Hab do not produce perturbations of the expansion or the density. Nor do perturbations of Eab and Hab produce rotational perturbations. 4. The Undisturbed Metric - -l"i\textbackslash{}cdot\textbackslash{}cdot\textbackslash{}cdot\textbackslash{}cdot--- Since in the unperturbed state the rotation and acceleration are zero, Ua must be hypersurface orthogonal. Ua...-s\textasciitilde{};o. ' where measure, the proper time along the world lines. As the surfaces 't = constant are homogeneous and isotropic they must be 3-surfaces of constant curvature. Therefore the metric can be written, where , d y 1 is the l1ne element of a space of zero or unit positive or negative curvature. 1.-.'c define t by, d. t I TI , clT :0. then \textbackslash{}cdot In this metric, . ' .. ' \textbackslash{}cdot\textbackslash{}cdot (prime denotes differentiatio.n- with respect to t )

\begin{figure}[h]
\centering
\includegraphics[width=0.9\linewidth]{images/page030_img01.png}
\caption{Extracted image from page 30}
\end{figure}

% --- Page 31 ---
\section*{Page 31}
\[
(5), (7) Then, by (20) \cdot\cdot \pm- G -:. - (f ft) \cdot3 1" ; \cdot ( 21 ) (2 fi. , If vve know the relation betvveen µ and 'vve may determine (l\_ \cdot % \Ve will consider the tv10 extreme cases, 1'"\. = 0 (dust) and }1 ~ (radiation). .Any physical situation should lie between these. 'b ;:: For 0 = By (20), rl.\ c onst \cdot \cdot\cdot -;, r1 \cdot :. 0 \cdot \cdot Q J f'(\ \cdot I 1.. D'2,- E \cdot = rt =: E const. \cdot \cdot M ) (a) For E 0 9 - E: n ' (cos t - 1 . h~r-M t -\cdot ) h~ - i E(~ t ; r- 1 ) 1t - ~ll'l -~ - - ~ 2 i ' o, (b) For T;' -- t:.1 n - f'I\ t3 -- M tz. - \cdot l 36 ' I '2. ) E represents the energy (kinet\cdotic + potential) per unit mass. If it is non-negative the universe will expand indefinitely, otherwtse it irvill eventually con tract again.
\]
\begin{figure}[h]
\centering
\includegraphics[width=0.9\linewidth]{images/page031_img01.png}
\caption{Extracted image from page 31}
\end{figure}

% --- Page 32 ---
\section*{Page 32}
\[
By tl1e Gauss Codazzi equations '"R; the curvature of the = bypersurface 'l const. is J EM - 2. - o~ E If )>O ) R - - 3 ~ M -= E \cdot ) [=O \cdot ) :> E <O ) b -3 /Y\ .- n2. ~ ,-:::: \cdot .c.: ) - ...t ---- - F'or = .?\ )J./"S ..,.\_\_.....,:/ \cdot ~2 - t - n \cdot\cdot 3 (2 -- ) . )l. '\.. - (2 I - 3 - E j-J\cdot -::- \cdot 0. ... \cdot ) \cdot fY\ o, (a) For > I~ - - .,.. R - (1 - -I ~.inh t 'L -- ' ( C.c\cdotS\cdoth t - I) -- 0(:,. 'l.. ,\cdot E ) E \cdot ) - o, (b) For E - ,..,.. - - t '. - - I 2. \cdot , t L 2. -<. - , ) ) <( 0 ] I - ( c) Fo1"' E 0 ~ - --' S . .i. YI t - ~ (c...ost -1) - E F .) ) 5. RotFttional Pert11rb at ions \_,,, \_\_ , \cdot\cdot\cdot--.-- \cdot- ... \cdot-\cdot.------.. \cdot------\cdot\_..-=-\cdot=ire-.-. By (6)
\]
\begin{figure}[h]
\centering
\includegraphics[width=0.9\linewidth]{images/page032_img01.png}
\caption{Extracted image from page 32}
\end{figure}

% --- Page 33 ---
\section*{Page 33}
\textbackslash{}cdot \textbackslash{}cdot w \textbackslash{}cdot\textbackslash{}cdot Q..b \textbackslash{}cdot For ) - - (\&),., (2 '2. 1 1ri.-:. For ) et w\textbackslash{}cdot -- - w e +-I - \textasciitilde{}) '+ ) f' - e --I w - '\} ) .. . - -- w Thus rotation dies away as the universe expands. This is in fact a statement of the conservation of angular momentum in an expanding un\textbackslash{}cdoti verse. 6. Perturbations of\_\_Densi.!il. fl\textbackslash{}cdot= For 0 we have the equations, - \textbackslash{}cdot - -t-8 fJ. \textbackslash{}cdot - -Te fJ - I 'a --rI: fA These involve no spatial derivatives. Thus the behaviour of one region is unaffected by tl1e behaviour of another. PeI'turbations ' will consist in some regions having slightly higl1er or lovrrer values of E than the average. If the untverse as wh "1..e has a value of E greater than zero, a small perturbation will till have E greater than n.nd will continue to expand. It will not cor.trsct to \textasciitilde{}ero

\begin{figure}[h]
\centering
\includegraphics[width=0.9\linewidth]{images/page033_img01.png}
\caption{Extracted image from page 33}
\end{figure}

% --- Page 34 ---
\section*{Page 34}
form a galaxy. If the universe has a value of E le8s than zero, a small perturbation can contract. However it will only begin contracting at a time 8 1" earlier than the wl1ole universe begins contracting, \textbackslash{}l\textbackslash{}There - -\textbackslash{}cdot .. is the time at vvl1ich the wl1ole m1iverse begins contracting. '(:" 0 = o. There is only any real instability vvhen E This case is of measure zero relative to all the possible values E can have . Ho-Yrever this cannot really be used as an arguement to dismiss it o. as there might be some reason why the universe should have E = = F\textbackslash{}cdotor a region with energy - cE , in a universe with E 0 \textbackslash{}cdot\textasciitilde{}E ( t'l.- t't r"l - t ... ) 4 12 >l - I T = se 1'1. ti E )'l.\textasciitilde{} ) ( f- ::: ----\textbackslash{}cdot '""t'" '1./\textasciitilde{} T \textbackslash{}cdot \textbackslash{}cdot \textbackslash{}cdot 2 \textasciitilde{}\textasciitilde{} For L: = O, 'l'3 . Thus the pe rturbation grows only as This is not fast enough to produce galaxies from statistical fluctuations even these i\textasciitilde{} could occt1r. Hovv0ver since an evolutionary universe has a partic1.e 9 9 horizon (Rindler(B) Penrosc( )) different parts do not communicate 9 in the early stages . This makes it ov0n more difficult for statistical fluctuations to occur over a rc;g: -.n until light had time to cross th0 region.

\begin{figure}[h]
\centering
\includegraphics[width=0.9\linewidth]{images/page034_img01.png}
\caption{Extracted image from page 34}
\end{figure}

% --- Page 35 ---
\section*{Page 35}
- - - n\textbackslash{}cdot \_I o i\_ ,\textbackslash{}cdot, ,,: .,,..., = - .., )A + 0 0 - V'- ... ' s a perturbation cannot contract unlesG it has a negative f .. bef'ore\textasciitilde{} value of E. The action of the pressure forces Iilake it still more difficult for it to contract . Eliminating 6, \textbackslash{}cdot\textbackslash{}cdot . . . I o.b .. \textbackslash{}cdot \textasciitilde{} 0..: -= + u l-l 0. t,..l <>-i b 1 0.. I.A to our approximation. I b :;:/ I OA.\textbackslash{}7 1 Ve I v b is tl1.e Laplacian in the hypersurface c-r = constant. o. \textbackslash{}cdot.\textasciitilde{}re represent the perturbatior1 as a sum of eigenfunctions S(n) of this operator, v1i1ere, ( ->) S <:. ,\textbackslash{}cdot cU. -;-.O s - - - Yt.7... \{f.o) Q'2. These eigenfunctions will be hyperspherical and pseudohyperspherical l'1armonics in cases (c) and (a) respectively anC plane \textbackslash{}vaves in case (b). In case (c) n \textbackslash{}ll!ill take only dj.screte \textbackslash{}cdot ,lues but in (a) and (b) it will take all positive values.

\begin{figure}[h]
\centering
\includegraphics[width=0.9\linewidth]{images/page035_img01.png}
\caption{Extracted image from page 35}
\end{figure}

% --- Page 36 ---
\section*{Page 36}
where µ is the undisturbed density \textbackslash{}cdot 0 \textbackslash{}cdot \textbackslash{}cdot \textbackslash{}cdot f'..'3 1 cng as fA '7 \textbackslash{}Vill gr0v1. 0 ' For >> f'o These perturbations grow for as long as light has not had time to travel a significant distance compared to the scale of the perturbation ( Q:. ). \textasciitilde{} Until that time pressure forces cannot act to even out perturbations. I/ B I I B rf. + 1.1/hen (11) (t1J ' 0 . ..,., ... e \{\_ .. t 1:3" r \textbackslash{}cdot !e obtain sound v-1aves whose amplitude decreases vii th time. T11ese results confirm those obtained by Lifsi1i tz and Khalatnikov(3). From the forgoing we see that galaxies cari.not form as the result; of the growth of small perturbations. ''Ve may x:pec t that other non- \textasciitilde{} gravitational force\textbackslash{}cdot will have an effect smaller than pressure equal

\begin{figure}[h]
\centering
\includegraphics[width=0.9\linewidth]{images/page036_img01.png}
\caption{Extracted image from page 36}
\end{figure}

% --- Page 37 ---
\section*{Page 37}
to one third of the density and so will not cause relative perturbations to grow faster than l: . To account for galaxies in an evolutionary uni verse we must assume tli.ere 11ver\textbackslash{}cdote finite, non-statistical, initial inhomogeneities. 7. Tpe Steady-stat\textasciitilde{}\_!.!D-jverse To obtain the st0ady - state universe \textbackslash{}Ve must add extra terms to 1 the energy-momentum tensor . Hoyle and Narlikar( 0) use, where, c 'CA ' J c \textbackslash{}cdot\textasciitilde{} 0.. .c. - ' J -r \textbackslash{}cdotb Sin.ce =- o 1 c..b' ) ( 21 ) c c h ,/,, \textbackslash{} - b \textasciitilde{}, ;d.. lJ\textbackslash{}cdot - Cl.+ I '-) b n CL - Q. b cl = o . (22) There is a difficulty here, if we require that the 11 Ci? field

\begin{figure}[h]
\centering
\includegraphics[width=0.9\linewidth]{images/page037_img01.png}
\caption{Extracted image from page 37}
\end{figure}

% --- Page 38 ---
\section*{Page 38}
should not produce acceleration or , in other words, that the matter created should have the same velocity as t11e matter already in existence.We must\textbackslash{}cdottJ;'=tr.> l1ave ..' .I (23) However since C is a scalar, this implies that the rotation of the medium is zero. On the other hand if (23) does not hold , the equations 11 (c.r. are indeterminate Raychaudhuri and Bannerjee( )) . In order to have a deterrninate set of equations we will adopt ( 23) but drop the requirement that Ca is the gradient of a scalar. The condition (23) is not very satisfactory but it is difficult to think of one more 12 satisfactory. Hoyle and Narlikar( ) seek to avoid this dj\_ff\textasciitilde{}iculty by taking a particle ratl1er than a fluid picture. Hovvever this has a 13 serious drawback since it leado to infinite fields (Hawking( )) . From ( 17), c c; e \textbackslash{}cdot ea-\textasciitilde{} \textasciitilde{} -= - 1\textasciitilde{}) - (.µ-1 1'l) \textbackslash{}cdot \textbackslash{}cdot 4 f \textbackslash{}cdot e = - 1 - \textbackslash{}. \textbackslash{}cdot

\begin{figure}[h]
\centering
\includegraphics[width=0.9\linewidth]{images/page038_img01.png}
\caption{Extracted image from page 38}
\end{figure}

% --- Page 39 ---
\section*{Page 39}
*ft f' For Tht1.s, small perturbations of density die avvny. Moreover equation ( 1 8) still 11olds, anc1 therefore rotational perturbations also die away. :t;q,uation (19) now becomes e e - r. \textbackslash{}cdot = - \textasciitilde{}I 'Z. 2I ( f-+ 3 ) + l .-----\textbackslash{}cdot 9 \_,, 1:- (1) :\textasciitilde{} \textasciitilde{} 3 ( 1 These results c onf j\_rm those obtained by Hoyle and l'iarlikar ( 4). Vie see therefore that galaxies cannot be formed in the steady-state uni verse by the gro,\textbackslash{}cdotrth of small perturb at ions. However\textbackslash{}cdot this does not exclude the possibility tha\textbackslash{}cdotc there migl"rt by a self-perpetuating system of finite perturbations which could produce galaxies . 1 16 (Sciama( 5), Roxburgh and Saff'man( )). ----------- 8. Gra\textbackslash{}cdot-v.. itational \textbackslash{}cdot1raves ... \textasciitilde{}- ----\textasciitilde{} \textbackslash{}cdot \textasciitilde{}Ve no\textbackslash{}1 consider perturba.tions of tl1e \textasciitilde{}Veyl tensor that do not arise from rotatj.onal or densi\textbackslash{}cdotcy p8rturbations, that is, \textbackslash{}cdotb -- LI b I =O f I 0. V<:. Multiplying ( 1 5) by lA c \textasciitilde{} and ( 16) by

\begin{figure}[h]
\centering
\includegraphics[width=0.9\linewidth]{images/page039_img01.png}
\caption{Extracted image from page 39}
\end{figure}

% --- Page 40 ---
\section*{Page 40}
we obtain , after a lot of reduction, - \textbackslash{}cdot\textbackslash{}cdot = r: 0 \textasciitilde{}b 17 In empty space with a non-expanding congruence Ua this reduces to the usual form of the l i nearised theory , The second term in (24) is the Laplacian in the hypersurface = "'r constant, acting on Eab \textbackslash{}cdot We will write Eab as a sum of eigensfunctions of this operator . . (1-1) where V = 0 I \} O\textasciitilde{}> v \_ yt\textasciitilde{} (1"1) Q b Q. '2. V 0 \textasciitilde{}o Cl'- \textbackslash{}cdot ) Then \textbackslash{}cdot \textbackslash{}cdot - \textbackslash{}cdot- \textbackslash{}cdot

\begin{figure}[h]
\centering
\includegraphics[width=0.9\linewidth]{images/page040_img01.png}
\caption{Extracted image from page 40}
\end{figure}

% --- Page 41 ---
\section*{Page 41}
Similarly , Then by (19) Substituting in (24) n'' Al"') [ n .... 3 (2 t D(\textasciitilde{}\} [ r2 (j-tt f'L) i 0 2(f fi-)] ::. 0 t i We may differentiate again and substitute for D' , For >>I \textbackslash{}'l J f2 and >/ t1 2. Cl - 1 so the gravitational field Eab decreases as and the "energy" r.... - 6 t(Eab Eab + Hab Hab) as 1 t.. \textbackslash{}cdot We might expect this as the Bianchi identities may be written, to the linear approximation, Therefore if the interaction with the matter could be neglected n D -1 cabcd would be proportional to and Eab ' Hab to \textbackslash{}cdot e In the steady-state universe when µ and have reached their . equilibrium values, R\textasciitilde{}\textasciitilde{} \textasciitilde{} /\textasciitilde{} fL \textasciitilde{} + i J \textbackslash{}cdot J \textasciitilde{} R La J - S \textasciitilde{} L\textasciicircum{}\{\textbackslash{}circ\} " 'R; \textbackslash{}cdot \textbackslash{}cdot o..b c c j b b :::: 0

\begin{figure}[h]
\centering
\includegraphics[width=0.9\linewidth]{images/page041_img01.png}
\caption{Extracted image from page 41}
\end{figure}

% --- Page 42 ---
\section*{Page 42}
c Thus the interaction of the 11 11 field vii th gravitational radiation is equal and opposi"i e to that of the matter . Ther'e is then no net n - 1 interaction, and Eab and decrease as \textbackslash{}cdot (E ,.,ab .1. The "energy" depends on se\textbackslash{}cdot ond derivatives ab-'-' 2 \textbackslash{}cdot of the me\textbackslash{}cdottr i Q. It is therefore proportional to the frequent.y squared times the energy as ineasured by the energy momentum pseudo- tensor, in a local co- moving Cartesian coordinate system .which depends only on first derivatives. Since the \textasciitilde{}requency will qe inversely proportional r2 , to the energy measured by the pseudo- tensor will be proportional - 4 ll to as for other rest mass zero fields . \textasciitilde{}L 9. of Gravitational AQ\_\textasciitilde{}orption lilf\textasciitilde{}es As vre have seen , gravitational waves are not absorbed by a perfect fluid .. Suppose however tl1ere is a small a1nount o1' viscosity. We may represent this by the addition of a term ). CJa ..b to the \textbackslash{}cdotA energy-momentum tensor, where is the coefficient of viscosity 1 ( El1le r s , ( 7) ) \textbackslash{}cdot Since we have (25) b he. (J ; -:::: 0 c.. b Ol. (26)

\begin{figure}[h]
\centering
\includegraphics[width=0.9\linewidth]{images/page042_img01.png}
\caption{Extracted image from page 42}
\end{figure}

% --- Page 43 ---
\section*{Page 43}
Equations (15) (16) become H \textasciitilde{};e = \textasciitilde{} (28) The extra terms on the right of equations (27), (28) are similar to conduction terms in Maxwell's equations will cause the wave to and e - "'t . decrease by a factor Neglecting expansion for the moment, \textasciitilde{} suppose we have a wave of the form , E \textasciitilde{}V"t E o.b ::::. - \textasciitilde{}be 0 . This will be absorbed in a characte:r.iatic time "-/). independent of frequency. By (25) the rate of gai\textbackslash{}cdot n of rest mass energy of the A- Ei. matter will be \textasciitilde{} A () which by ( 19) will be ?. v -2. . Thus the 2. 0 Fz .-2 .:>.:. )) \textbackslash{}cdot This confirms that the available energy in the wave is 4density of available energy of gravitational radiation will decrease as (l -4 in an expanding universe . From this we see that gravitational radiation behaves in much the same way as other radiation fields. In the early stages of an evolutionary universe vvhen the temperature was ver\textbackslash{}cdoty high we migl1t expect an equilibrium to bo sei up between black-body electromagnetic diation and black- body \textasciitilde{} gravitational radition. Since they both have \textbackslash{}cdotciivo polarisations their '

\begin{figure}[h]
\centering
\includegraphics[width=0.9\linewidth]{images/page043_img01.png}
\caption{Extracted image from page 43}
\end{figure}

% --- Page 44 ---
\section*{Page 44}
energy densities should be equal. .As the uni verse expanded they wouJ. ':. both cool adiabatically at the same rate. As we know the temperatu\textasciitilde{}c of black-body extragalactic elec1 radiation is less than \textasciitilde{}omagnetic 5\textasciicircum{}\{\textbackslash{}circ\}:< , the terrperc.:.ture of the blac1c-body gravitational radiation n1us'..; be also less \textbackslash{}cdotthan this which l be absolutely undetectable. Now vvou\textasciitilde{}. the energy c \textbackslash{}cdotf gravitational does not contribute to the radi\textasciitilde{}\textbackslash{}cdot :\textasciitilde{}ion ordinary e11ergy momentum tensor TI I'reverthele ss it will have an ... ab \textbackslash{}cdot e.ctive gravi tatio11al ef'fect . By the expansion equation, \textbackslash{}cdot For incoherent gravitational radiation at frequency v , E'l. 2. -i \textasciitilde{} ::. 0 ')) 2. -2. E v But the energy density of the radiation is 4 0 \textbackslash{}cdot e I 2 I -:; - 5 8 - Z \}J- (f 3 \}1.) \textbackslash{}cdotf \textbackslash{}cdot G - \{ t- \textbackslash{}vhe1'3 µG is the gravitational 1tenergyn density. Thus gravitational radiation has an active attractive gravitational effect. It is interesting t11at tl1is seems to be just half that of electromagnetic radiat ior1. 1 8 1 It 11as been suggested by I-Iogarth ( ) and Hoyle and Na'Y'likar ( O), that there may be a connection between the absorption of radiation and the Arrow of Time. Thus in universes like the steady-state, in which all electromagnetic radiation emitted is :ventually absorbed by other matter, the Absorber theory would predic retarded solutions of

\begin{figure}[h]
\centering
\includegraphics[width=0.9\linewidth]{images/page044_img01.png}
\caption{Extracted image from page 44}
\end{figure}

% --- Page 45 ---
\section*{Page 45}
the Maxwell equations while in evolutionary universes in which electroinagnetic radiation is not completely absorbed it would predict advanced solutions. Similarly, if one accepted this theory, one would expect retarded solutions of the Einstein equations if and only if all gravitational r\textbackslash{}cdotadiation emitted is eventually absorbed by other ma.t ter. Clearly this is so for the steady - state universe since /\textbackslash{} will be constant. In evolutionary universes will be a function of time. r V1e ',ii/ill obtain complete absorption if >t o1 .r diverges. Now f (" '.' a ga:-;, " 1 2 \textasciitilde{} oc. T2 \textbackslash{}'\textbackslash{}There T is the temperature. For a monatomic gas , T fl - , d:. therefore the integral will diverge (just). However the expression used for viscosity assumed that the mean free path of the atoms was small compared to the scale of the disturbance. Since the mean free 1 n 1 path oc µ- oe -3 and the wavelength (2 - ' the me 8.n free path will 0::. eventually be greater than the wavelength and so the effective viscosity () -1 \textbackslash{}cdot will decrease more rapidly than 1 L Thus there will not be conwlcto absorption and the theory would not predict retarded solutions. However this is slightly academic since gravitational radiation has no\textbackslash{}cdot::; yet been let alone invest.igated to see vvhether it corresponds d\textasciitilde{}tectcd, to a retarded or solution. adva.\textasciitilde{}ced

\begin{figure}[h]
\centering
\includegraphics[width=0.9\linewidth]{images/page045_img01.png}
\caption{Extracted image from page 45}
\end{figure}

% --- Page 46 ---
\section*{Page 46}
R-ef-erences (1) Boilllor W.B . ; M. N., :L1.Z, 104 (1957) . (2) r.ifshit z E. M.; J.Phys . u.s . s . R., 10 , 116 (1946). (3) Lifshitz E. M. , Khalatnikov I . N. ; Adv. in Phys . , 12 , 185 (1963) . (iq6'f) (4) Irvine V"J. ; Annals of Physics . \textasciitilde{} J \textasciitilde{}\textbackslash{}cdotz 2. (5) Jordan P. , Ehlers J ., Kundt W. ; Ahb . Akad . Wiss . Mainz ., No . 7 (1960) \textbackslash{}cdot . ( 6) Kundt \textbackslash{}V. , Trllinper M. ; Ahb. Akad . Wiss . Mainz . , No . 1 2 ( 1 964)., (7) TrUmper M. ; Contributions to Actual Problems in Gen . Rel . (8) Rindler W.; 116 , 662 (1956) . M.r\textasciitilde{}., (9) Penrose R. ; Article in ' Relativity Topology and Groups ', Gordon .and Br each ( 1 964) . (10) Hoyle F . , Narlikar J . v.; Proc. Roy . Soc ., A, 'fil, 1 (1964) . (11) Raychaudhuri A., Banerjee S.; Z. Astrophys., .2§., 187 (1964) . (12) Hoyle F ., Narlikar J . V. ; Proc.Roy. Soc ., A, 184 (1964) . \textasciitilde{}8.f., 8. vv. ; (13) Hawking Proc . Roy. Soc ., A, ... :in prj,nt. \textasciitilde{} \textasciitilde{}/3> (/\textasciitilde{}.£) (14) Hoyle F. , Narlikar J . V. ; Proc . Roy . Soc ., A, 273 , 1 (196 3 ) (15) Sciama M. N., 1J5-., 3 (1955) . D . \textasciitilde{}v. ; (16) Roxburgh I.'1V., Sattman P. G. ; 1'11. N. , J.12 181 (1965). (17) Ehlers J . ; Ahb.Akad.Wiss . Mainz.Noll . (1961 ). (18) Hogarth J.E. ; Proc . Roy . Soc . , A, 267 , 365 (1962) .

\begin{figure}[h]
\centering
\includegraphics[width=0.9\linewidth]{images/page046_img01.png}
\caption{Extracted image from page 46}
\end{figure}

% --- Page 47 ---
\section*{Page 47}
3 CHl\textbackslash{}.Pl'En Gravitational .Liadi\&tion In An Expanding Unive1\textbackslash{}cdotse Gru.vi tational radiation in e:npty asymptotically flat space has been examined by means of aoymptotic exoansions 1 4 by a nL.mber of \&uthors . ( - ) \textbackslash{}cdot'.Chey find that the different components of "che out,soing radiation field "peel of'f that 11 , is , they so different powers of the affine raJial distance . vS If one v1ishes to inves\textbackslash{}cdottigc:.:.te hoi'l this beha\textbackslash{}cdotriour is :nodi .i:-iPd by the of matter , one is faced with a pr\textasciitilde{}sence difficult\textasciitilde{}r that doAs not arise in the case o.r, say , ele\textasciitilde{}tromagnetic radiation in matter . For this one can consid.er the radiation trc.Lvellint; thro11gh infinite t1nii'orm medium that is static a11 apart :from "t;he disturbance cr,3ated by the ra\textbackslash{}.i.iat;ion. In the case oi' cro.vi tcttional t;hi:J ir;; no\textbackslash{}cdott pos,\_;ible . racli\textasciitilde{}tion ' .!!Or, if che were initially static, its own self me\textasciitilde{}ium 5ravitatio\textasciitilde{} 1r1ouia. (;ause it to contr a ct in on itself and it v1ould cease to be static . Hence one is forced to investigate gravi\textbackslash{}cdottatio\textbackslash{}cdot\textasciitilde{}al in matter that is either contracting or expandin;. raQi\textasciitilde{}tivD in Jhapter 2 , we identify the Weyl or conformal i\textasciitilde{}s 1 \textbackslash{}cdottensor with the freA gravitational field and the 1\textasciitilde{}icci-t ensor Rctb with the contribution of thA matter to the curvature. Instead of considerin; gravitational radiation in '

\begin{figure}[h]
\centering
\includegraphics[width=0.9\linewidth]{images/page047_img01.png}
\caption{Extracted image from page 47}
\end{figure}

% --- Page 48 ---
\section*{Page 48}
acymptotically flat space , tna.t is, space that \textasciitilde{}pproach\textasciitilde{}s flat sc; .. .t.ce at large radial distances , \textbackslash{}ve consider it in rts,11n:pt o tica.lly conformally flat space . .1\textbackslash{}.s it i.;.; only conformally flat , the and the density of \textasciitilde{}icci-tensor m\textasciitilde{}tter need not be zero . :'o avoid essentiall;yy non- gravitational phenon1ena. sue h as sound wavgs , we will conbiner gravitationalra\textasciitilde{}iation travellin6 through dust . It was shown in Chapter 2 that a conformally flat universe filled v1ith dust must have one of the metrics: I L l ri 1 f4 f (a) \textbackslash{} Ls - i .: S-2. i ( cl \{; 4 - cl.f i - s\textbackslash{}cdot ,,\textbackslash{} 2 ( c D l s. ' r. :t C/Jl c 7' ).) 't' ; A ( c) (;. c as: i \textbackslash{}cdot- 1.) . \textbackslash{}cdot1 :i I\textasciitilde{} \textasciitilde{} l. ). J2 cl I '2. 'l. t: . C( ;.. (b) \textasciitilde{} \_\textbackslash{})\_ ;. A t \textasciitilde{} (I, 2 ) 2. - L c S' 'l - (c) \_\_J <j\_ - - fype(a) represents a universe in which the mutter expands from \textbackslash{}cdotthe initial sinJularity with insu fficient ener3y to reach infinity and so falls back to another a\textasciitilde{}aj.n si11gularity. It i s therefore unsuitable for a dj.scussion of \textbackslash{}cdot

\begin{figure}[h]
\centering
\includegraphics[width=0.9\linewidth]{images/page048_img01.png}
\caption{Extracted image from page 48}
\end{figure}

% --- Page 49 ---
\section*{Page 49}
ravi\textbackslash{}cdottationril radiation b;y- a method of .::..sympto\textbackslash{}cdottic expansions 0 si11ce one cu.11no\textbackslash{}cdott <\_;et an inf ini ve C...:.istun(\textasciitilde{}e from -eh\textasciitilde{} source. Type (b) is the \textasciitilde{}instein-De Sitter universs in \textasciitilde{}hich ::he mat\textbackslash{}cdotter has just sufficient energy to reach infinity . It is thus a )ec i o.l case. D. Norman ( 5) has investit;<J..ted the S 1 "peeling off" behaviour in this case usi11c; Penrose ' s co11f ormal 6 technique ( ). He was however forced to make certain assumpt - ior1s about the movement of the matter .\textbackslash{}cdot1hich vrill be sho\textasciitilde{}\textbackslash{}cdot!n to 1 be false . ;:oreover , he \textbackslash{}vas misled by the soec i al nc.ture o""' the \textasciitilde{}instein-De Sitter universe in which affine ard lumir:osi .. y dis\textbackslash{}cdottance:\_, differ . Another reason for not considerin;:; ra.::li<:...tion in the t\textbackslash{}cdot:instein-De Sitter uni verse is that it is unstable. fhe passage of a gravitational wave will cause it to contract a\textasciitilde{}ai\textbackslash{}cdoti e ventually and develop a s in.;ularity . de will therefore consider radiation in a universe o! t, pe ( c) \textasciitilde{}1hich corresponds to the general case vJhere tn.e 1 rna\textbackslash{}cdottteI' i.3 expanding th more th.: n enough ener6y to avoid \textbackslash{}\textbackslash{}cdotJi contructin;; again . 2. fhe Newman-Penrose Formalism \textasciitilde{}\textbackslash{}cdotJe employ the notation of l\{ewman and Penrose. (3) A r-- l nf-) rn f'-. tetrad of null vectors , f-/. ('(\} / t'I 1\textbackslash{}cdot L is introduced ,

\begin{figure}[h]
\centering
\includegraphics[width=0.9\linewidth]{images/page049_img01.png}
\caption{Extracted image from page 49}
\end{figure}

% --- Page 50 ---
\section*{Page 50}
\[
nf' (Y)i~ L,,., (/,r\ M 1iil1v.: lf- -- (lt' /l?rl; [ , v1nere : lf' p \cdot- f.J. f- . Mr -} M ;; l f l) (L = vve label vectors with a tetrad th~se p.- ('L: M ru) fl~ P1'~ 2 = ().\_ tetrad irdices are raised \lith vhe notric an~ lo~ered - - c~b uq -- -- ? 0 l 0 0 ( ?\_\_:\_\_!) ( 0 0 0 \cdot 0 0 0 -1 l o -I 0 0 L pv O.b 2v 2p. have - vie - 7 -~ ) . ().- h m v L rr1 \cdot- fA fl l f'L11 + p- ,~ v \_\_ fV1 fYl (2 j\_) .. II \cdot- - IA - Ricci rotation coefficients are defined by : 3) ( 2 -
\]
\begin{figure}[h]
\centering
\includegraphics[width=0.9\linewidth]{images/page050_img01.png}
\caption{Extracted image from page 50}
\end{figure}

% --- Page 51 ---
\section*{Page 51}
In f'ac\textbackslash{}cdott it is more convenient to in of 1J-;.,'el 1e ':.:01:l\{ t \textbackslash{}cdot\textasciitilde{}rms 1 complex combinations of coefficients defined as ro\textasciitilde{}ation f ollO\textbackslash{}'/S : /\textbackslash{} - ( - I \textasciitilde{} I -y - \textbackslash{}cdot 2. lt .L /Al. l I n v\textasciitilde{} 1Yl 1iVif - v)\textbackslash{} - 'I I \textbackslash{}cdot 'I - - ( l - I - 1 \{v<.; v I L //.\_\_,.;..; \textbackslash{}cdot .'.:l )'l I 3\textasciitilde{}1 v fA - - 0 - Lr: I - - v 1'l!J fl1 - \textbackslash{}cdot \textbackslash{}cdot- .... \textbackslash{}cdot \textbackslash{}cdot- - \textbackslash{}cdot \textbackslash{}cdot - . - - -

\begin{figure}[h]
\centering
\includegraphics[width=0.9\linewidth]{images/page051_img01.png}
\caption{Extracted image from page 51}
\end{figure}

% --- Page 52 ---
\section*{Page 52}
\[
3 . Coordi.:1atcs --- .Lil ee i:\cdotl e \fJrnan a nd Penrose , introduce a nu ll coo1:-dir1ate \Je u..( = xj) i) lL \cdot v <~ . ~ o I r L L vte take ~ \cdot Thus ':Till be 3eodesi c and irrotati onal . This implies v -- 0 (l - ( f - ~ cc- ::: - £ :L ~ + vie take be parallelly ~o transporu~d along IJ.'b.is gives \cdot - s) -~ \cdot 1~s a se concl c oordina\cdotte \':e talce an affine parameter l f" long the t:;eodesics 1\_ r:\cdot f L ~ I (-:? f" Lt ) 1 X X ~' and arc coordinates that label the geo6.e...:.ic t\\cdot10 i n the surfv .ce L.l = cons t . \cdotu. L': f.i. 0 l Y't ,,v, - l). X : L = I\ ,' i1hu s 1 (1. b) In these { f) ( ~ l =
\]
\begin{figure}[h]
\centering
\includegraphics[width=0.9\linewidth]{images/page052_img01.png}
\caption{Extracted image from page 52}
\end{figure}

% --- Page 53 ---
\section*{Page 53}
\[
The lield ~~u~tions Je 1nay calculate the Ric ci and ;1Je.1l tensor com ion8n't;s f r om 1,1 1 the rela\cdottions a..\_b c..-.:.L o-. b c: (L {.-{ 't \cdot- ~ "") c, ' - ,- - \cdot J ~ ") ..\_L b Using thA combinations of rotation coefficients defined <ind with /"C -= TT = £ = 0 v1e have p () ( J. 10) j) ~ I!) G' - \cdot\cdot\cdot- ( \cdots.. ]) 'f i2) /)d. (1.1.3) - ..,\_ j)~ ( 1 I~) - Pr ( :s . /~) y ,,\ -- -- (s. I b) Pr (~\cdot- tJ) vv {l . IY)
\]
\begin{figure}[h]
\centering
\includegraphics[width=0.9\linewidth]{images/page053_img01.png}
\caption{Extracted image from page 53}
\end{figure}

% --- Page 54 ---
\section*{Page 54}
\[
-- f , po ' f l C\) ~\cdot ~ ~, {? ~ J ('5 . c;- -- ( r ;\_ ) -t ( . - b' \cdot- r 'l c'/ ) - -' y; 6 f ->. - f" J<X ~ r~ ~ O(~ ~? ~" j\_ - G \cdot- 7 ,- \cdott- /\ r '2::\_ 1) - f f - fe ) J~ f- ..,. (;} - ~ ~) ). - r;/t,' , ( oi. .,. 3 1 T J - 4 r - ?J.2 tr-2vf\cdotir'f'?.\_,~IT fv - ( (? f 1J ~ J (S fA - J f ) " f .,- (I. , I - l. !.. ") b -r T i ^{\circ}I" ). 0 - \cdot r - s cr ;- i \cdot- f - --r; - ). )\ ,\_ 11 f ) 2""' ~:, t). = l. "'L - ,\,,-. - r ^{\circ}' - i l J'' ( ft \_ (J\_ \_ L~ f) = v - - ). L J
\]
\begin{figure}[h]
\centering
\includegraphics[width=0.9\linewidth]{images/page054_img01.png}
\caption{Extracted image from page 54}
\end{figure}
\end{document}
